\chapter{Introduction}


\section{Project Overview and Rationale}

Indoor positioning systems have become increasingly critical in modern applications where Global Positioning System (GPS) signals are unavailable or unreliable. From warehouse asset tracking and hospital patient monitoring to augmented reality navigation and smart building automation, the demand for accurate, cost-effective indoor localization solutions continues to grow. While GPS provides excellent outdoor positioning accuracy, its signals cannot penetrate building structures effectively, creating a significant gap in location-based services for indoor environments. Historically, technologies such as Wi-Fi and Bluetooth using Received Signal Strength Indicator (RSSI) measurements were deployed to address this challenge. However, RSSI is highly susceptible to severe multipath fading, environmental noise, and physical obstructions, making reliable positioning exceedingly difficult without dense anchor networks and exhaustive calibration. Conversely, Ultra-Wideband (UWB) technology offers remarkable centimeter-level precision but commands a prohibitive infrastructure cost, demanding strict clock synchronization and expensive multi-anchor setups, which effectively limit its scalability for ubiquitous deployment.\par

This project presents the design and implementation of an advanced, low-cost Indoor Positioning System (IPS) that fundamentally shifts away from conventional RSSI and expensive UWB paradigms. Instead, it leverages the cutting-edge LD2450 24GHz Millimeter-Wave (mmWave) radar sensor orchestrated by an ESP32 microcontroller. The LD2450 provides innate, high-precision detection and continuous tracking of up to three moving targets simultaneously within its field of view, actively computing the spatial X and Y coordinates without requiring computationally intensive external algorithms. The ESP32 microcontroller, featuring integrated Wi-Fi capabilities, a dual-core processor, and scalable architecture, operates as an efficient host controller to parse the mmWave UART data and facilitate wide-area network telemetry. This synergy delivers highly reliable, privacy-preserving tracking—an immense advantage over optical cameras—at a fraction of the cost associated with UWB arrays.\par

The fundamental principle underlying this newly proposed system relies on the Frequency Modulated Continuous Wave (FMCW) radar mechanics of the LD2450. Rather than estimating distance through signal degradation (as with RSSI), FMCW radar actively transmits a continuous mmWave signal varying in frequency and analyzes the corresponding reflection from the target. By calculating the time of flight and angle of arrival from the reflected waves, the sensor precisely determines both range and angular position. This enables robust tracking that remains largely unfazed by ambient lighting, varying room temperatures, or minor non-metallic obstructions. The system directly interfaces with these Cartesian outputs via serial communication protocols, translating complex electromagnetic wave physics into easily actionable positional data vectors.\par

\subsection{Low-Cost Reproducible Testbed Approach}

A distinguishing feature of this work is its emphasis on reproducibility, accessibility, and robust validation through a highly affordable hardware platform. The system is designed with real-world applicability and seamless scalability as its core development principles. To thoroughly evaluate the LD2450 sensor's native capabilities against theoretical tracking limitations, a standardized experimental testbed was established. By positioning the sensor on a controlled axis and documenting target traversal coordinates over pre-measured paths, the experiment ensures rigorously consistent testing conditions. The utilization of readily available ESP32 microcontrollers significantly streamlines the ecosystem. Open-source programming environments and standard serial parsing logic enable rapid replication and extension by other researchers and academic practitioners with genuinely minimal upfront investment. This approach comprehensively democratizes high-fidelity indoor positioning research by systematically removing the financial barriers typically associated with high-end proprietary hardware, intricate time-of-flight synchronization matrices, and restrictive closed-source licensing models.\par

\subsection{Three-Phase Architectural Evolution}

A key innovation of this work is the structural implementation plan, featuring a three-phase architectural evolution designed to balance accuracy validations, spatial scalability, and practical implementation complexity. The first phase executes a \textit{Single-Node Validation} system where a standalone LD2450 sensor paired with an ESP32 host successfully tracks subjects in a bounded indoor environment. This stage focuses heavily on strictly parsing the serial byte protocol emitted by the sensor, filtering extraneous noise reflections, and mapping the innate X and Y values against real-world ground-truth scenarios to establish a reliable baseline precision. The second phase systematically extends the paradigm into a \textit{Multi-Node Spatial Coverage} configuration. By bridging multiple LD2450/ESP32 clustered nodes via ubiquitous Wi-Fi or ESP-NOW networks to a localized aggregator, the architecture elegantly resolves intrinsic line-of-sight constraints and occlusion vulnerabilities. When a target passes behind an opaque pillar or exits the primary sensor's field of view, the interconnected secondary nodes seamlessly continue the tracking operation. The third phase completely realizes a \textit{Real-Time Visualization and System Integration} ecosystem. During this stage, the parsed positional metrics accumulated from the distributed multi-node network are transmitted to a centralized server, rendering the moving targets dynamically upon a digital layout. This demonstrates full-stack readiness for end-user deployment.\par

The cohesive use of standardized JSON data packet structures and networked aggregation mechanisms ensures backwards compatibility across all three phases, allowing smooth, incremental deployment and comprehensive diagnostic testing. This methodical, phased approach not only radically simplifies the iterative development and debugging cycles but also explicitly provides clear migration pathways for diverse application requirements. Each technological phase reliably builds upon its predecessor, maintaining rigid compatibility with spatial coordinate mapping and asynchronous data streaming protocols. This ensures that the low-level UART interfacing and noise-filtering algorithms refined during the preliminary single-node testing remain fundamentally integral and profoundly valuable as the system scales into a complex, multi-room tracking environment.\par

By consciously pivoting away from conventional techniques, this system wholly eliminates the necessity for empirical path-loss calibration protocols and speculative machine learning regressions. Instead of relying on predictive neural networks to smooth erratic signal strengths, the system trusts the robust, physics-based Digital Signal Processing (DSP) natively executing within the LD2450 radar module. This drastically reduces the overarching computational overhead demanded from the ESP32 host, allocating processor resources toward efficient network routing and multi-sensor fusion rather than complex positional inference math.\par

\section{Objectives}

The primary objectives of this project are strictly tailored to the integration and deployment of mmWave sensors:

\begin{enumerate}
    \item \textbf{Protocol Interfacing and Data Extraction:} Design and successfully implement a robust UART parsing algorithm on the ESP32 to synchronously capture, decode, and structure the continuous serial data stream generated by the LD2450 mmWave sensor, isolating precise target coordinates.
    
    \item \textbf{Performance and Feasibility Analysis:} Systematically evaluate the performance, precision, and cost-to-benefit ratio of the 24GHz mmWave technology specifically compared to previously explored, expensive UWB alternatives and highly volatile RSSI methodologies.
    
    \item \textbf{Architecture Development:} Implement the three progressive system phases: (a) Single-Node Validation for establishing baseline accuracy bounds, (b) Multi-Node Spatial Coverage for combating environmental occlusion, and (c) Real-Time Visualization for practical end-user tracking demonstration.
    
    \item \textbf{Accuracy and Range Validation:} Conduct methodical physical testing to achieve and empirically validate the sensor's advertised positioning precision and maximum effective tracking range within practical indoor settings, rigorously documenting any limitations concerning angular field-of-view constraints.
    
    \item \textbf{Network Aggregation and Synchronization:} Develop reliable wireless transmission logic that enables multiple dispersed ESP32 nodes to funnel their respective mmWave tracking metrics to a central processing terminal without suffering from data packet collisions or severe latency degradation.
\end{enumerate}

\section{Contributions}

This project makes several crucial contributions to modern indoor positioning system infrastructure and deployment strategies:

\begin{enumerate}
    \item \textbf{Cost-Effective UWB Alternative:} Demonstrates a practical, privacy-preserving tracking methodology utilizing the LD2450 radar that achieves highly reliable multi-target spatial tracking at a fraction of the hardware cost traditionally commanded by UWB infrastructure.
    
    \item \textbf{Occlusion-Resistant Architecture:} A scalable multi-node network topology specifically engineered to counteract line-of-sight physical occlusions by strategically overlapping sensor detection zones, ensuring unbroken surveillance of mobile entities.
    
    \item \textbf{Standardized Processing Pipeline:} The creation of a well-defined, modular software pipeline that translates raw, byte-level radar outputs into standardized spatial coordinates, facilitating immediate integration with broader Internet of Things (IoT) ecosystems and higher-level graphical user interfaces.
    
    \item \textbf{Simplified Deployment Paradigm:} By completely eliminating the extensive environmental calibration previously required for RSSI path-loss models and discarding computationally heavy machine learning regressions, this methodology presents a genuinely plug-and-play positioning solution.
    
    \item \textbf{Real-Time Distributed Telemetry:} Engineered communication protocols that confidently aggregate localized data vectors from distributed sensor nodes while maintaining exceptionally low latency, enabling true real-time location mapping essential for critical health-monitoring or high-speed automation applications.
    
    \item \textbf{Open-Source Hardware Validation:} Comprehensive empirical validation of the commercially available LD2450 module operating in tandem with the ESP32, providing future researchers with verified, actionable performance metrics and pre-configured integration templates for accelerated academic or industrial replication.
\end{enumerate}

\section{Organization of the Chapters}

The remainder of this report is logically organized to reflect the progression of the project as follows:\par

\textbf{Chapter 2: Literature Review} comprehensively surveys existing research paradigms in indoor positioning, critically examining the inherent limitations of RSSI and the fiscal disadvantages of UWB. It formally introduces the theoretical foundation of Frequency Modulated Continuous Wave (FMCW) millimeter-wave radar and highlights recent advances in LD2450-related literature.\par

\textbf{Chapter 3: Problem Identification} explicitly articulates the driving technical challenges that motivated the abandonment of RSSI techniques. It details the destructive nature of multipath fading, the immense burden of environmental calibration, and the strict constraints regarding line-of-sight blockages, ultimately validating the strategic pivot toward mmWave networks.\par

\textbf{Chapter 4: Proposed Methodology} presents the exhaustive system design blueprint. This encompasses the nuanced three-phase evolution architecture, detailed schematic layouts showing ESP32 to LD2450 micro-connections, comprehensive flowcharts dissecting the UART serial parsing logic, and the topology configuring the expanded multi-node wireless aggregator.\par

\textbf{Chapter 5: Experimental Results and Discussion} meticulously documents the physical testing protocols and corresponding spatial tracking accuracy. It analyzes observed data regarding sensor precision, boundary detection limitations, multi-node overlap performance, and thoroughly discusses the comparative findings against the discarded positioning techniques.\par

\textbf{Chapter 6: Conclusion and Future Work} neatly summarizes the definitive achievements of the redefined project. It reflects upon the profound success of migrating to mmWave technology, outlines unresolved edge cases discovered during tracking testing, and explores ambitious future enhancements such as deploying three-dimensional multi-floor tracking or integrating complex activity-classification behavioral analytics.
